\label{sec:3pDuoram}
\textit{DUORAM} is an \textit{DORAM} implementation by Vadapalli et al. \cite{vadapalli2023duoram}.
It can either be used as 3-party or 2-party protocol.
This research will focus solely on the 3-party variant.
\textit{DUORAM} uses (2,2)-DPFs to represent \textit{Private Information Retrieval (PIR)} and PIR-Writing queries (\ref{sub:PrivateInformationRetreival}).
All the following protocols and notations are as described by Vadapalli et al. \cite{vadapalli2023duoram}.

At a high level, the database $D$ is additively secret shared among two non-colluding parties $P_0,P_1$.
Both parties hold the additive shares of the target index $i^*=i_0+i_1$.
This index $i^*$ is represented by the shifting vector $e_i = e_{i_0+i_1}$ representing an evaluated DPF:
\[
\mathbf{e}_i = 
\begin{cases}
    1, & \text{if} \quad j = i, \\
    0, & \text{else}.
\end{cases}
\]

For writing accesses both parties additionally hold additive shares $M_0,M_1$ of the target value $M=M_0+M_1$.

Each party holds three RDPFs in form of \textit{RDPF-Triples}.
Evaluating them, either by using \code{EVALFULL} (\ref{sub:TraversingTheDPFTree}) or while generating both RDPF-Triples on target index $i^*$ (\ref{sub:BGIDPFConstruction}), returns six pairs of word vectors $(t_b^k,v_b^k), b\in\{0,1\}\text{ and }k\in\{1,2,3\}$ with the property that $t_0^1+t_1^1=t_0^2+t_1^2=t_0^3+t_1^3=e_{i^*}$ and $v_0^1+v_1^1=v_0^2+v_1^2=v_0^3+v_1^3=M\cdot e_{i^*}$.

Also, each party carries a \textit{Blinding Factor ($\zeta$)} that will be used to mask its share of the database $D_i$.
The blinded share of party $P_b, b\in\{0,1\}$ is being denoted by $\widetilde{D}_b\leftarrow D_b+\zeta_b$.
Each party carries the blinded database of the other party.
This means party $P_b, b\in\{0,1\}$ holds $\widetilde{D}_{1-b}$.

This 3-party variant also has a stateful auxiliary part $P_2$.
It holds neither the database $D$ nor index shares, and merely facilitates the computation.
This allows executing (2+1)-MPC protocols.
This means, three parties compute the result on two secret inputs with one party facilitating the computation without holding secret inputs itself.
3-Party Duoram uses (2+1)-MPC protocols to replace the trusted source \cite{vadapalli2023duoram}.

\subsection{Reading and Writing}
\label{sub:DuoramReadingAndWriting}
\paragraph{Reading:}
\label{par:DuoramReading}
The client wants to read at index $i^*$ from $D$.
This notation is equivalent to $D[i^*]=\langle D,e_{i^*}\rangle$.
The algorithm works as follows:
\begin{enumerate}
    \item $P_2$ selects $p\in\{0,1\}^w$ uniformly at random and sends $y_0\leftarrow-\langle \zeta_0,t_1^3\rangle + p$ to $P_0$ and $y_1\leftarrow\langle\zeta_1,t_0^2\rangle -p$ to $P_1$.
    \item $P_0$ outputs $\text{read}_0\leftarrow\langle D_0+\widetilde{D}_1, t_0^1\rangle - \langle\zeta_0,t_0^3-t_0^1\rangle + y_0$.
    \item $P_1$ outputs $\text{read}_1\leftarrow\langle D_1+\widetilde{D}_0,t_1^1\rangle - \langle\zeta_1,t_1^2-t_1^1\rangle +y_1$.
\end{enumerate}

\paragraph{Writing:}
\label{par:DuoramWriting}
Both parties want to obtain shares of a new Database $D'$, whose $i^*$-th value is updated by $M$.
This is equivalent to updating the read value $r=D[i^*]$ by $M-r$ or completely overwriting $r$ by $M$ at position $i$.
Recall the word-vectors $v_b^k$ with $v_0^k+v_1^k=e_{i^*}\cdot M$.
This implies, parties $P_0,P_1$ locally set $D'_b\leftarrow D_b+v_b^1\text{ for }b\in\{0,1\}$.
Since writing is a combination of reading and updating, one needs two RDPF-Triples for executing a single write operation.

\paragraph{Refresh Blinds:}
\label{par:RefreshBlinds}
The blinded, exchanged shares become "stale" after a single update operation.
This means $\widetilde{D}'_b\neq D_{b-1}+\zeta_{b-1}, \text{ for } b\in\{0,1\}$ after executing a write operation.
One needs to refresh them to guarantee correctness in the following operations.
This means party $P_b, b\in\{0,1\}$ has to refresh $\widetilde{D}_b$ so it matches $D_{1-b}+\zeta_{1-b}$ and party $P_{1-b}$ has to refresh its blind $\zeta_{1-b}$ before.

This specific operation is not of great importance in this research.
If not specified differently, $\text{RefreshBlinds}$ will just be adapted to the given datasize.
The benevolent reader may read the more detailed research by Vadapalli et al. \cite{vadapalli2023duoram}.

\subsection{Preprocessing- and Online-Phase}
\label{sub:DuoramPreAndOnlinePhase}
\textit{DUORAM} works in two phases, namely the \textit{Preprocessing Phase} and the \textit{Online Phase}.
The intention is to shift the calculation and communication expensive part to the preprocessing phase to have everything ready for reading and writing.

\subsubsection*{Preprocessing-Phase}
\label{subsub:DuoramPreprocessing}
Preprocessing generates RDPF-Triples (\ref{sub:DistributedPointFunctions}).
Depending on the parameters given, one can also generate distributed comparison functions (\ref{sub:DistributedComparisonFunctions}), multiplication triples, and addition triples (\ref{sub:BeaverTriples}).
Each RDPF-Triple is generated on the same share, thus all three DPFs by $P_0$ have the same index, and all DPFs held by $P_1$ share the same index too.
Of course, these do not have to be equal.

\textit{DUORAM} uses a (2+1)-Party MPC protocol for generating their DPFs, which allows to replace the trusted source by a stateless third party $P_2$.
Both parties $P_0,P_1$ choose a random index $\widetilde{ri}_0,\widetilde{ri}_1$ with $\widetilde{ri}_0\oplus\widetilde{ri}_1=ri$.
The DPFs will be generated on these target indices using the BGI-construction (\ref{sub:BGIDPFConstruction}).
This results in DPF-Triples $(\overline{k}_0^1,\overline{k}_0^2,\overline{k}_0^3)$ and $(\overline{k}_1^1,\overline{k}_1^2,\overline{k}_1^3)$.
$P_0$ will then send $\overline{k}_0^2$ to $P_2$ and $P_1$ will send $\overline{k}_1^3$ to $P_2$.
Note the usage of XOR shares here while the rest of \textit{DUORAM} works on additive shares.
Therefore, using a MPC share conversion procedure, the XOR shares are turned to additive shares $ri_0,ri_1$ with $ri_0+ri_1=ri$.
This also affects the word vectors $v_b^i,t_b^i$ representing the leafs and flag bits.
These will be turned to additive shares too.
A more detailed explanation of the MPC protocol to change from XOR to additive shares can be found in the research of Vadapalli et al. \cite{vadapalli2023duoram}.

Preprocessing results in both parties $P_0,P_1$ holding DPF-Triples and their evaluations in form of additive shares. 
$P_2$ holds the DPFs $(\overline{k}_0^2, \overline{k}_1^3)$.

\subsubsection*{Online-Phase}
\label{subsub:DuoramOnlinePhase}
This phase includes read and write operations executed by $P_0,P_1$.
Both parties hold DPF-Triples on random target-indices $ri_0,ri_1$.
To write or read from the non-random target index $i^*$ both parties need to execute the \textit{Cyclic-Shift Protocol} (\ref{sub:CyclicShiftingProtocol}).
\textit{DUORAM} uses cyclic shifting in two flavors to postpone decision that have not been made in the \textit{Preprocessing-Phase}:
Namely postponing the decision of where to read/write or to postpone the decision of what to write.

\paragraph{Postponing the decision of \textit{where} to read or write:}
This protocol adjusts the shares of the shifting-vector $e_{ri}$ to the target index $e_{i^*}$.
This changes the target index from random $ri=ri_0+ri_1$ to the new target index $i^*=i^*_0+i^*_1$.
Now, both parties need to obtain vectors $v_0+v_1=e_{i^*}$ to represent the shifting vector.

They exchange $i^*_b-ri_b$, reconstruct $cs=(i_0^*+i_1^*-ri_0-ri_1)$ and cyclic right shift $v_b$ by $cs$.
$P_2$ receives $i^*_b-ri_b$ from $P_b$ and can also compute $cs$.

\paragraph{Postponing the decision of \textit{what} to write:}
Recall the goal of the writing-operation: Adding $M=M_0+M_1$ at the already shifted target index $i^*$ to the database $D$.
At the moment, evaluating the RDPF at position $i^*$ results in a random word.
Therefore, one must set the final correction word $F$ to change the random result into the message $M$.

$P_0,P_1$ exchange $M_b+F_b$ (with $F$ being the final correction word).
This allows both parties to reconstruct $M-(v_0[i^*]+v_1[i^*])$ which will serve as the new final correction word $F'$ to write $M$ in the desired location.
The cyclic shift results in the shifted vectors $\tilde{t}_b^i,\tilde{v}_b^i, (i\in\{1,2,3\},b\in\{0,1\})$ with the final properties described in \ref{sec:3pDuoram}.
Finally, all three parties execute the writing or reading protocol from \ref{sub:DuoramReadingAndWriting}.

\subsection{Low-Level Reading and Writing} \label{sub:AdvancedReadingAndWriting}
Having established a theoretical foundation of the \textit{DUORAM} construction in \ref{sec:3pDuoram}, we will now dive further into reading and writing.
All the mentioned protocols in this subsection are taken directly from Vadapalli et al. \cite{vadapalli2023duoram}.

\subsubsection*{Reading}\label{subsub:advancedReading}
Recall the fundamentals of the 3-Party Duoram-Reading-Protocol from section \ref{par:DuoramReading}.
Parties $P_0,P_1$ want to read from target index $i^*$ in the shared database $D$.
This is equivalent to $D[i^*]$.
$P_0,P_1$ are called \textit{peers} in the following.
Party $P_2$ is called the \textit{server} in the following.
It facilitates the communication and does not hold any unique information on its own.

This means $P_0$ holds $(D_0,\xi_0,\widetilde{D}_1,i^*_0)$, $P_1$ holds $(D_1, \xi_1,\widetilde{D}_0,i^*_1)$, and $P_2$ holds $(\xi_0,\xi_1)$.
Recall $\xi_b$ being the blinding factor of party $P_b \text{ with }b\in\{0,1\}$ defined by $\widetilde{D}_b\leftarrow D_b\oplus\xi_b$. 

$P_0,P_1$ return $\text{read}_0,\text{read}_1$, such that $\text{read}_0+\text{read}_1=D[i_0^*+i_1^*]$. 

\paragraph{Preprocessing}\label{par:advancedReadingPreproc}
Advanced preprocessing does not extend the already introduced preprocessing phase (\ref{subsub:DuoramPreprocessing}).
It, too, generates two XOR-shared RDPF-Triples $(\overline{k}_b^1,\overline{k}_b^2,\overline{k}_b^3) \text{ with } b\in\{0,1\}$.
The distribution of triples between parties also stays the same.
Fully evaluating both triples would each return three XOR-shared vectors $\overline{v}_b^k,\overline{t}_b^k,F_b^k \text{ with }b\in\{0,1\},k\in\{1,2,3\}$.
Note $F$ not being overlined as it stays XOR-shared and will not be converted to an additive share.
Recall $t$ being the bit-vector holding the flagbit of leaf $i$ at position $i$, $v$ being the vector holding all leaf labels and $F$ being the final-correction-word (\ref{sub:BGIDPFConstruction}). %//TODO nochmal Final Correction Word nachschauen und bei bedarf weiter ausführen
$t_b^k$ has the property of $e_{i^*}=t_0^1+t_1^1=t_0^2+t_1^2=t_0^3+t_1^3$ with $e_{i^*}$ being the shifting vector representing a unit vector holding a single $1$ at index $i^*$. 

Due to these properties, advanced reading does not require the RDPF-Triples to be fully evaluated but only needs the six flagbit vectors $t_b^k$. 
As evaluation on XOR-shared RDPF-Triples returns XOR-shared vectors, these need to be converted to additive shares using an (2+1)-MPC protocol (\ref{sec:3pDuoram}).
This provides consistency throughout the implementation as the database itself is additively shared and therefore does not need to be converted to XOR-shares.

\paragraph{Online-Phase}\label{par:advancedReadingOnline}
The online phase combines both cyclic shifting (\ref{sub:CyclicShiftingProtocol}) and the already introduced reading protocol (\ref{par:DuoramReading}).
This means reading first executes cyclic shifting to decide where to read and then executes the basic reading protocol. 
It results in parties $P_0,P_1$ returning $\text{read}_0,\text{read}_1$ with $\text{read}_0+\text{read}_1=D[i^*_0+i^*_1]$
Figure 5 in the research of Vadapalli et al. \cite{vadapalli2023duoram} provide a nice overview over the reading online protocol.

\subsubsection*{Writing}\label{subsub:advancedWriting}
Recall the fundamentals of the writing operation from section \ref{par:DuoramWriting}.
It takes three operations to complete a single writing operation: reading, updating, and refreshing the blinds.
Vadapalli et al. combine both updating and refreshing in a single \code{update} operation.
This requires a single read operation to be executed in advance, thus consuming two separate RDPF-Triples - one for reading and one for updating.
These triples can be generated in two separate preprocessing phases executed directly before each online-phase or alternatively in a single preprocessing phase delivering multiple RDPF-Triples.

This means, after executing a single read, $P_0$ holds $(D_0,\xi_0,\widetilde{D}_1,i_0^*,M_0)$, $P_1$ holds $(D_1,\xi_1,\widetilde{D}_0,i^*_1,M_1)$, and $P_2$ holds $(\xi_0,\xi_1)$.
After updating, which includes the refresh of the blinds, $P_0$ holds $(D'_0,\xi'_0,\widetilde{D}'_1)$ and $P_1$ holds $(D'_1,\xi'_1,\widetilde{D}'_0)$ such that $\widetilde{D}'_0=D'_0+\xi'_0$, $\widetilde{D}'_1=D'_1+\xi'_1$, and $D'_0+D'_1=D_0+D_1+(M_0+M_1)*e_{i^*_0+i^*_1}$.
$P_2$ gets $(\xi_0,\xi_1)$ with the same properties as above.
In short, message $M$ is written at position $i^*$ into the shared database $D$ resulting in the new database $D'$ which, still, is shared additively among the peers.

\paragraph{Preprocessing}\label{par:advancedWritingPreproc}
Advanced preprocessing, too, does not extend the already defined preprocessing from \ref{subsub:DuoramPreprocessing}.
This means $P_0,P_1$ again generate the same RDPF-Triples on random indices and distribute them as above. 
However, more parameters are extracted compared to the reading operation (\ref{par:advancedReadingPreproc}).

During RDPF generation $P_0,P_1$ learn XOR-shares of the random final correction words (\ref{sub:BGIDPFConstruction}) $F_0^k,F_1^k\text{ with }k\in\{1,2,3\}$.
These will be important during the online phase as they incorporate the message $M$.
Note, that both the message $M$ and its additive shares $M_0,M_1$ are the same for both RDPF-Triples.

Finally, $\widetilde{v}_b^k,\widetilde{t}_b^k$ are converted to additive shares.
Note that despite the final corrections words being XOR-shared they are not be converted to additive shares as the RDPF itself stays XOR-shared.

\paragraph{Online-Phase}\label{par:advancedWritingOnline}
The online phase for writing combines cyclic shifting, updating the database, and refreshing the blinds.
This means cyclic shifting needs to determine both where and what to write.
Therefore, $P_0,P_1$ exchange $(i^*_b+ri_b),(M_b+F_b^k)$.
$P_2$ will receive all of that but $(M_0+F_0^1),(M_1+F_1^1)$
These values allow the peers to reconstruct three final correction words $F^k\leftarrow(M_0+F_0^k+M_1+F_1^k)$ and the server and the peers to reconstruct the offset $S\leftarrow(i_0^*+ri_0+i_1^*+ri_1)$.
The server, too, reconstructs $F^k$ but only for $k\in\{2,3\}$.
Note that $F_b^k$ is the negative sum of the according RDPF's leaf labels $F_b^k=-\sum^n_i{v_b^k[i]}$.

Both peers update their share of the database $D_b$ for all indices $i$ as follows:
$$
D_b'[i]\leftarrow D_b[i]+(v_b^1[i-S]+F^1\cdot t_b^1[i-S])
$$
Combining these values for both peers at the target index $i^*$ will result in:
\begin{flalign*}
   D_0[i^*]+D_1[i^*]+(v_0^1[i^*]+t_0^1[i^*]\cdot F^1)+(v_1^1[i^*]+t_1^1[i^*]\cdot F^1) \\
    = D[i^*]+v^1[i^*]+F^1\cdot(t_0^1[i^*]+t_1^1[i^*]) \\
    = D[i^*]+v^1[i^*]+(M-(v_0^1[i^*]+v_1^1[i^*])) \\
    = D[i^*]+v^1[i^*]-v^1[i^*]+M \\
    = D[i^*]+M \\
    = D'[i^*] 
\end{flalign*}

Recall $t_0^k[i]+t_1^k[i]=e_{i^*}[i]=0$ for every $i\neq i^*$.
Due to
\begin{align*}
    D_0[i]+D_1[i]+(v_0^1[i]+t_0^1[i]\cdot F^1)+(v_1^1[i]+t_1^1[i]\cdot F^1) \\
    = D[i] + v_0^1[i]+v_1^1[i] 
    = D + v^1[i] 
    = D[i] + (e_{i^*}[i]\cdot M) \\
    = D[i]
\end{align*}
$D[i]$ stays at the initial unchanged value for every index $i\neq i^*$.

$P_2$ is not directly involved while updating. 
It solely updates its own copy of the peers' blinds for data-consistency in the next read operation.

Figure 6 in the research of Vadapalli et al. \cite{vadapalli2023duoram} provides a great overview over the writing online-phase.